% Generated by docs/user-manual/tables/return_period_flow.py -- do not hand-edit.
% Source: all forecast days of releases 2026-11-06 to 2026-11-13.
\begin{table}[htbp]
\centering
\small
\begin{tabular}{l r r r r}
\toprule
Return period & Forecasts & Min flow & Median flow & Under 10 m\textsuperscript{3}/s \\
\midrule
\multicolumn{5}{l}{\textbf{GEOGLOWS}} \\
  2-year & 45,325 & 5.0 & 52.8 & 16.8\% \\
  5-year & 3,164 & 5.0 & 59.5 & 19.1\% \\
  10-year & 1,257 & 5.0 & 48.3 & 24.6\% \\
  25-year & 421 & 5.0 & 23.1 & 36.6\% \\
  50-year & 248 & 5.0 & 10.5 & 48.4\% \\
  100-year & 1,875 & 5.0 & 6.6 & 74.3\% \\
\midrule
\multicolumn{5}{l}{\textbf{Flood Hub}} \\
  2-year & 415 & 17.6 & 916.9 & 0.0\% \\
  5-year & 105 & 51.9 & 248.5 & 0.0\% \\
  20-year & 15 & 79.3 & 247.0 & 0.0\% \\
\bottomrule
\end{tabular}
\caption[Discharge against return period]{%
  \textbf{Discharge against return period.}
  Every forecast day of 8 releases (2026-11-06 to
  2026-11-13), pooled. Read down the median column: if the climatology
  behind a return period is sound, a rarer event carries \emph{more} water.
  Where the median falls as the return period rises, the return period is not
  measuring what it claims --- the same defect as the blackout described above,
  at a scale too thin to see on the map. GloFAS is absent because it publishes
  no discharge, only ensemble exceedance counts. GEOGLOWS figures exclude the blacked-out region in 7 of 8 releases, so the defect shown here is the part of it that survives outside that box.}
\label{tab:return-period-flow}
\end{table}
